\section{Introduction}\label{sec:introduction}

Images of resolved galaxies encode the projected distribution of their
stellar populations, dust, and substructure. Following contours of constant
surface brightness reduces this information to radial profiles of intensity,
ellipticity, position angle, centroid, and departures from pure ellipses while
retaining important signatures of disks, bars, twists, lopsidedness, and tidal
structure. This approach grew from the scanning of photographic plates with
microdensitometers \citep{Fraser1967}. \citet{Carter1978} introduced a Fourier
description of deviations from fitted ellipses, and the iterative
intensity-harmonic algorithm of \citet{Jedrzejewski1987} subsequently became
the standard basis for the IRAF \texttt{Ellipse} task. Higher-order terms were
soon used as physical diagnostics of boxy and disky structure
\citep{Bender1988,Bender1989}. More recently, \citet{Ciambur2015} introduced
eccentric-anomaly sampling and the \texttt{Isofit} and \texttt{Cmodel} tasks
for strongly flattened or non-elliptical systems. The classical algorithm is
now available in Python through \texttt{photutils.isophote}
\citep{Bradley2025}, while AutoProf provides a modern automated pipeline aimed
at deep survey imaging \citep{Stone2021}.

The continued value of one-dimensional isophotal analysis lies partly in its
balance between fidelity and efficiency. A single S\'ersic model is compact
and interpretable but cannot describe the full diversity of galaxy structure;
a many-component decomposition or highly flexible basis can approach the
image more closely, but at greater computational cost and with more parameter
degeneracy \citep{Peng2010,Erwin2015,Robotham2017}. A sequence of flexible
isophotes occupies a useful middle ground: it makes fewer assumptions about
the radial light distribution, preserves structural changes with radius, and
remains straightforward to inspect across large samples. Deeper imaging
strengthens this use case. Faint stellar haloes and intracluster light often
require additional model components, while their inferred structure becomes
increasingly sensitive to sky subtraction, scattered light, masks, and the
extended point-spread function \citep{Abraham2014,Duc2015,Trujillo2016}.
Isophotal profiles remain actively used in this regime: HSC imaging has traced
the stellar haloes and ellipticity profiles of thousands of individual massive
galaxies to 100 kpc \citep{Huang2018}; SGA-2020 provides surface-brightness
profiles for 383,620 nearby galaxies from the DESI Legacy Imaging Surveys
\citep{Moustakas2023}; and semimajor-axis profiles have been used to measure
brightest cluster galaxies and intracluster light near
$30\,\mathrm{mag\,arcsec^{-2}}$ \citep{Kluge2020}. Thus, one-dimensional
analysis is complementary to two-dimensional modeling rather than displaced
by it, and its favorable fidelity--efficiency balance is especially valuable
for deep, multi-band surveys.

Current implementations nevertheless leave several practical limitations.
Efficiency is a central issue because the iterative algorithm repeatedly
samples both the trial isophote and the local radial gradient at many radii;
the benchmarks presented in this work show that this cost can be reduced
substantially without sacrificing the recovered profiles. A second limitation
of the classical IRAF/\texttt{photutils} formulation is that its harmonic fits
use ordinary least squares and do not accept the per-pixel variance or RMS maps
now routinely produced by imaging pipelines. Equal weighting discards useful
information about spatially varying noise and limits rigorous uncertainty
propagation. The low-surface-brightness regime adds further difficulty because
uncertain radial gradients can amplify small residuals into large geometry
corrections, particularly in the presence of imperfect masks, nearby sources,
or background structure. Existing software addresses different parts of
these problems: IRAF remains an important reference but its legacy environment
is difficult to integrate into contemporary survey pipelines \citep{Tody1986};
\texttt{photutils.isophote} provides an accessible implementation of the
classical method; and AutoProf emphasizes automated, stabilized surface
photometry \citep{Stone2021}. Capabilities such as eccentric-anomaly sampling,
variance-aware fitting, faint-outskirt controls, multi-band analysis, and
reproducible diagnostics therefore remain scattered across different tools
and workflows.

We present \isoster\ (ISOphote on STERoid), an open-source Python package
designed to integrate these requirements while retaining the transparent
Jedrzejewski--Ciambur harmonic framework. \isoster\ accelerates the calculation
through vectorized path sampling, compiled numerical kernels, and selective
reuse of radial-gradient measurements, making repeated isophotal analysis
practical within large imaging-survey pipelines. It supports weighted
least-squares fitting from per-pixel variance maps, classical and
eccentric-anomaly sampling, higher-order harmonic analysis, and explicit
controls for unstable low-surface-brightness geometry. Reproducible
single-band and multi-band workflows, including an experimental joint-fit
mode, two-dimensional model reconstruction, validated configuration, status
information, serialization, and common quality-assurance diagnostics are
provided through one interface. In this paper, we describe the implementation
and evaluate its computational performance, scientific fidelity, depth, and
robustness against established isophotal tools using controlled mock galaxies
and representative survey images. Our aim is not to replace every existing
package, but to provide a fast and flexible implementation that allows
one-dimensional isophotal analysis to serve as a routine measurement layer for
modern imaging surveys.

